% Source: https://artofproblemsolving.com/community/c3h1724670_numbering_problem_from_2011_test
% Stable identity: aops-topic-1724670-q01
\textbf{Upper bound.} Give each number on the list its smallest prime factor. These primes are all different, since two numbers sharing it would have a common prime factor. Every number $n$ with $20<n<50$ has smallest prime factor $p$ with $p^2\le n$ or $p=n$ (if $n$ is composite, $p\le\sqrt{n}<\sqrt{50}$, so $p\in\{2,3,5,7\}$). So the smallest prime factors come from $\{2,3,5,7\}\cup\{23,29,31,37,41,43,47\}$, which has $4+7=11$ elements; the primes $11,13,17,19$ cannot occur. So the list has at most $11$ numbers.

\textbf{Construction.} The list $25,27,32,49,23,29,31,37,41,43,47$ works: $32=2^5$, $27=3^3$, $25=5^2$, $49=7^2$ and the seven primes pairwise share no prime factor.

The answer is $\boxed{11}$.
